% 功能概述：统一 camera-ready 结果表的视觉样式，复用早期 IEEE 主表的深蓝表头及四色方法分组。
% 主要职责：PaperTableSetup 设置局部字号和间距；PaperHead 设置白色粗体表头；PaperGroup 和 PaperNote 分别排版分组及说明。
% 代码库关系：由 main.tex 在导言区载入，供 latest_results、native_results、historical_results 的编号表调用。
% 输入输出与副作用：只影响显式调用的表格；颜色无数值排名含义，不改动指标、样本或统计口径。
\usepackage{tabularx}
\usepackage{caption}
\captionsetup[table]{font=small,labelfont=bf,skip=6pt,justification=centering}
\definecolor{HdrBg}{HTML}{1E3A5F}
\definecolor{RetBg}{HTML}{DBEAFE}
\definecolor{DirBg}{HTML}{DCFCE7}
\definecolor{KnwBg}{HTML}{FEF3C7}
\definecolor{AgtBg}{HTML}{FCE7F3}
\definecolor{SecTxt}{HTML}{1E293B}
\newcommand{\PaperTableSetup}[1]{%
  #1\renewcommand{\arraystretch}{1.38}%
  \setlength{\tabcolsep}{5pt}%
  \setlength{\aboverulesep}{0pt}%
  \setlength{\belowrulesep}{0pt}%
  \arrayrulecolor{HdrBg}}
\newcommand{\PaperHead}[1]{{\color{white}\bfseries #1}}
\newcommand{\PaperGroup}[3]{%
  \rowcolor{#2}\multicolumn{#1}{l}{\color{SecTxt}\bfseries\itshape #3}\\}
\newcommand{\PaperNote}[1]{%
  \par\vspace{4pt}\begin{minipage}{\linewidth}%
  \footnotesize\raggedright #1\end{minipage}}
